% =====================================================================
\section{Lessons}
\label{sec:lessons}
% =====================================================================

\paragraph{When a joint controller is worth building}
For planners of this kind, joint search earns its cost only when one
interval may move a tenant far relative to its size \emph{and} the layers
share a resource. The rate-limited lattice moves a tenant by at most
one-eighth to one-third of its replica cap per interval; the wide lattice
by half or more. Kubernetes' default HPA scale-up policy, which may double
a deployment or add four pods every 15\,s~\cite{k8shpa}, resembles the wide
lattice at 6 to 16 replicas, so the regime in which coordination pays is
not exotic; but in our plant it was also the regime in which the planner
did worst. Whether a deployment is in the separable regime can be checked
before anything is built: probe the disagreement between joint and per-knob
best responses on recorded states, as Fig.~\ref{fig:width}(b) does, in the
spirit of Bristol's interaction measure~\cite{bristol1966}. Below a few
percent disagreement, three well-tuned loops will do.

\paragraph{What to build instead}
In this plant the bill is model spend, and the decision that moves it is
the tier. A per-layer stack with a tuned, capped autoscaler and the
cheapest adequate model matched the joint controller on the matrix and beat
it on unseen demand, also when held to the cluster caps, and on the
live-fitted plant. What the joint
controller added that the stack lacked was a hard per-tenant budget, which
it kept by refusing service. A budget is an operator policy and should be
enforced, and reported, as one: a gateway that refuses AI calls once a
tenant's budget is spent makes the trade explicit, whereas an objective that
prices the refusal through a capped violation term hides it. The final
campaign gave a per-layer stack exactly that gate. On the synthetic matrix
the gate rarely fired and the gated stack tied the joint controller; on real
BurstGPT demand, where the budgets bind, the joint controller kept them at a
lower cost than the gate did. That is the one place our evidence favours a
budget-aware planner over a reactive layer, and it is narrow: arms that
simply overspent the budget scored better still.

\paragraph{Fixes interact with the actuation model}
The published controller's switching penalty ratchets replicas upward, which
the audit flagged as a defect, and the corrected penalty was the treatment's
best variant on the instant-actuation tuning slice. Under one-interval lag
the ranking reversed: the ratchet held replicas the plant could not deliver
in time anyway, while the corrected controller paid for every replica that
arrived late, and it lost to the reactive stacks it had tied. Hysteresis in a
controller is a bet on the actuation model, and a variant selected under one
model needs a registered test under the others before it is called a fix.

\paragraph{How to evaluate a joint controller}
Every comparator defect we found favoured the joint controller.
Table~\ref{tab:checklist} turns them into rules, each paired with the
defect it caught here and the effect of correcting it. None of the rules is
new on its own~\cite{chameleon,hoefler2015}; what this study adds is a
measurement of how much each was worth in one realistic evaluation, and
the observation that they compound: the first two corrections alone
removed a 27.5\% lead, and the fourth removed a 60\% one.

\begin{table}[t]
\caption{Rules for comparing a joint controller with per-layer
controllers, the defect each caught in this study, and the effect of
correcting it on the joint controller's measured advantage ($\Delta J$ as
a percentage of the comparator's $J$; negative: joint better).}
\label{tab:checklist}
\centering\tablefont
\begin{tabular}{@{}>{\raggedright\arraybackslash}p{0.25\textwidth}
                >{\raggedright\arraybackslash}p{0.36\textwidth}
                >{\raggedright\arraybackslash}p{0.31\textwidth}@{}}
\toprule
Rule & Defect it caught & Effect (campaign) \\
\midrule
One cost model for all arms & Only baselines paid a $1.4581\times$
  inference surcharge; their cache was pinned at 128\,MB &
  $-27.5\%\rightarrow-0.4\%$ (EVICTION\_PARITY) \\
Interior tuning optima, per cluster size & HPA's optimum on a grid edge &
  $-0.9\%\rightarrow0.0\%$ (RT-H1) \\
Same caps, budgets and action set for all arms & Re-tuned HPA and the
  cheapest-tier stack ran up to twice the small cluster's replica cap; only
  the joint controller kept budgets or could shed &
  capped HPA still a tie (JN-B4); capped stack: small-cluster loss
  $+0.0732\rightarrow+0.0066$ (exploratory); with caps, a budget gate and an
  outage price for all, no joint win on the matrix (LEVEL) \\
A cheapest-adequate per-layer stack & Registered tier rule ranked tiers by
  latency, at 3.4$\times$ the cost & $-60.4\%\rightarrow+4.7\%$ (JN-B1) \\
A decoupled-search ablation & Coordination never isolated &
  1.6\% of $J$, inside the margin (JN-C1) \\
Count requests served; price outages & $J$ priced a shed request only
  through a capped violation term & 10.1\% of AI requests unserved; priced
  and level, no win on the matrix (SA, LEVEL) \\
Held-out demand & Every workload periodic, comparators tuned on them &
  $+23.3\%$ against the cheapest-tier stack (HO-B1); $+6.5\%$ with the stack
  capped (exploratory) \\
\bottomrule
\end{tabular}
\end{table}

\paragraph{For multi-loop self-adaptive systems}
The question of when separate adaptation loops need a coordinator is
usually argued from the loops' interactions~\cite{weyns2013,kephart2007}.
Our results suggest it can also be answered by measurement, and that the
answer depends on two properties an engineer controls: how far each loop
may move per decision, and whether the loops' effects meet in a shared
resource. Rate limits, which operators impose for stability anyway, make
the joint problem locally separable and the coordinator nearly redundant.
This is one plant and one family of planners; whether the regularity holds
for other multi-loop systems is an open, testable question.
